\begin{idstatement}{C.2.VII.3}{Proved sign for class VII normalization}
\label{st:C.2.VII.3}
Let $f:Y\to B$ be a proper connected-fibre holomorphic
map from a smooth compact complex threefold with
$b_2(Y)=0$ to a smooth curve, whose general fibre is a
complex two-torus. Let $mS$ be a full fibre with
irreducible reduced nonnormal support, and suppose
that the minimal resolution $Z$ of its normalization
is class VII. Then
\[
 e(S)\ge0.
\tag{C.VII.7}
\]
No restriction on the singularity type or conductor
multiplicity of $S$ is made.
\end{idstatement}
\begin{proof}
The canonical evaluation and adjunction comparison
(A.CD.1) give $\omega_S^{\otimes m}=\OO_S$.
The divisor sequence and Hirzebruch--Riemann--Roch
give $\chi(\OO_S)=0$, because $c_1(\OO_Y(S))=0$
rationally. These arguments require only $b_2(Y)=0$;
the integral condition and Picard-variation assumption
in the later part of A.5.12 are not needed here.
Let $\eta:T\to S$, $\pi:Z\to T$, and
$D>0$ be the maps and divisor of A.5.3. Put
$C=\Gamma_{\mathrm{red}}$ and
$A=(\eta^{-1}C)_{\mathrm{red}}\subset T$.
Thus $K_Z+D$ is torsion, and C.2.VII.1--C.2.VII.2
apply.

We compute $e(T\setminus A)$ through the actual
resolution. As in C.2.VII.2, an exceptional curve
$E$ with coefficient zero in $D$ has
$K_Z\cdot E=D\cdot E=0$.
Adjunction and negative self-intersection then give
$p_a(E)=0$ and $E^2=-2$.
These curves are smooth rational and disjoint
from $D_{\mathrm{red}}$.
Their negative-definite intersection matrix makes
each connected configuration a tree: a repeated
edge or a cycle would give a nonnegative square
for a nonzero vector in the matrix. The intersections
are transverse, since their local intersection
multiplicities are one; three curves cannot meet
at a point, since that would give a triangle.
Connectedness of resolution fibres shows that a
fibre either consists entirely of these curves or
lies entirely in $D_{\mathrm{red}}$.
The zero-coefficient fibres lie outside $A$, since
a fibre over $A$ must meet the strict conductor.

Let $a$ be the total number of zero-coefficient
exceptional curves, and let $b$ be the number of
exceptional fibres lying in $D_{\mathrm{red}}$
whose images in $T$ lie outside $A$. On
$Z\setminus D_{\mathrm{red}}$ the former fibres
are still present and the latter are absent.
A tree with $a_j$ rational curves has Euler
number $a_j+1$, so replacing it by its image
point subtracts $a_j$. Adding back the image
point of an absent positive-coefficient fibre
adds one. Away from these finitely many fibres,
$\pi$ is the required isomorphism of complements.
Additivity of compactly supported Euler
characteristic therefore gives the exact equality
\[
 \begin{aligned}
 e(T\setminus A)
   &=e(Z)-e(D_{\mathrm{red}})-a+b\\
   &=b_2(Z)-r(D_{\mathrm{red}})-a+b
   \ \ge0 .
 \end{aligned}
\tag{C.VII.8}
\]
The inequality follows because the rational curves
counted by $r(D_{\mathrm{red}})$ and by $a$ are
distinct, and their total is at most $b_2(Z)$ by
C.2.VII.1.
Finally $\eta$ is a biholomorphism
$T\setminus A\to S\setminus C$. Hence
\[
 e(S)=e(T\setminus A)+e(C)\ge0
\]
by C.2.VII.2. This proves the assertion for the
original reduced support and all full multiplicities.
\end{proof}


% C.2.VII.4 from part_C2_classVII.tex

